\section*{Referencias}
\addcontentsline{toc}{section}{Referencias}

\begin{itemize}[leftmargin=1.27cm, itemindent=-1.27cm, label={}, itemsep=12pt, parsep=0pt]

\item Asamblea Nacional de Gobiernos Regionales. (2024, 18 de junio). ANGR preocupada por estancamiento de la inversión en la Red Dorsal. \textit{Gestión}. \url{https://gestion.pe/economia/angr-preocupada-por-estancamiento-de-la-inversion-en-la-red-dorsal-noticia/}

\item Channel News Perú. (2023, 22 de noviembre). \textit{La Red Dorsal: Un modelo fallido frente a la urgencia de conectividad digital}. \url{https://channelnewsperu.com/la-red-dorsal-un-modelo-fallido-frente-a-la-urgencia-de-conectividad-digital/}

\item Cuarto Poder. (2021, 23 de mayo). \textit{Red Dorsal Nacional: La obra de fibra óptica que costó millones y casi no se usa} [Video]. YouTube. \url{https://www.youtube.com/watch?v=3fbD6yyRpRI}

\item Red Dorsal Nacional de Fibra Óptica. (2017). \textit{Scribd}. \url{https://es.scribd.com/document/358263533/red-dorsal-nacional-de-fibra-optica-docx}

\item Rojas, J. (2023, 29 de octubre). La Red Dorsal Nacional de Fibra Óptica opera solo al 8\% Las deficiencias que no permitieron la conectividad de la población. \textit{Infobae}. \url{https://www.infobae.com/peru/2023/10/29/la-red-dorsal-nacional-de-fibra-optica-opera-solo-el-8-las-deficiencias-que-no-permitieron-la-conectividad-de-la-poblacion/}

\item Tesis UNI. (2020, 15 de septiembre). \textit{Evaluación del proyecto Red Dorsal Nacional de Fibra Óptica en el Perú} [Video]. YouTube. \url{https://www.youtube.com/watch?v=Chyz-MJvF0M}

\item Universidad Nacional Federico Villarreal. (2022). \textit{Análisis e impacto socioeconómico de la Red Dorsal Nacional de Fibra Óptica} [Tesis de licenciatura, Universidad Nacional Federico Villarreal]. Repositorio Institucional UNFV. \url{https://repositorio.unfv.edu.pe/server/api/core/bitstreams/2ddcbcb5-277e-412b-b12e-03cd52de7db0/content}

\end{itemize}